\noindent
{\Large \textbf{\section{Результаты работы}}} \par

Измерения выполнены на компьютере с процессором Apple M4 под управлением macOS. Привязка потока к ядру в macOS недоступна, поэтому планировщик мог выбрать любое из ядер. Вывод программы разбит на три таблицы по числу экспериментов и приведен в Листингах~\ref{lst:result-line}, \ref{lst:result-size}, \ref{lst:result-assoc}.



\newpage
\noindent{\Large \textbf{\subsection{Размер строки}}} \par
\begin{lstlisting}[language=,caption=Определение размера строки,label=lst:result-line]
1. Размер строки
шаг, байт   нс/обращение
        8      0.294
       16      0.338
       32      0.991
       64      1.158
      128      9.324
      256     14.172
      512     29.137
Размер строки: 128 байт
\end{lstlisting}
\par
При шагах 8 и 16 байт время обращения составляет 0.29-0.34 нс, при шагах 32 и 64 байта - около 1 нс. При шаге 128 байт время возрастает до 9.3 нс, то есть в 8.1 раза по сравнению с шагом 64 байта. Это наибольший рост между соседними шагами, поэтому размер строки кэша составляет \textbf{128 байт}.

\vspace{8pt}

\noindent{\Large \textbf{\subsection{Размер кэша}}} \par
\begin{lstlisting}[language=,caption=Определение размера кэша,label=lst:result-size]
2. Размер кэша
размер, КБ   нс/обращение
         4      0.917
         8      0.921
        16      0.920
        32      0.921
        64      0.922
       128      0.923
       256      5.260
       512      5.143
      1024      5.303
      2048      5.403
      4096      6.580
      8192      7.803
     16384     20.447
     32768     85.423
     65536     92.886
    131072     96.719
    262144     97.692
\end{lstlisting}
\par
\begin{enumerate}
    \item Для массивов от 4 до 128 КБ время обращения постоянно и составляет около 0.92 нс. При размере 256 КБ оно возрастает до 5.3 нс. Следовательно, размер кэша первого уровня составляет \textbf{128 КБ}.
    \item Для массивов от 256 КБ до 8 МБ время держится в диапазоне 5.1-7.8 нс. При размере 16 МБ оно составляет 20.4 нс, а при 32 МБ - 85.4 нс. Для массивов от 64 МБ время выходит на значение порядка 90-98 нс. Следовательно, размер кэша второго уровня составляет \textbf{16 МБ}.
    \item Третий уровень кэша в замерах не выявлен: после второго скачка новых ступеней нет.
\end{enumerate}

\vspace{8pt}

\newpage
\noindent{\Large \textbf{\subsection{Ассоциативность}}} \par
\begin{lstlisting}[language=,caption=Определение ассоциативности,label=lst:result-assoc]
3. Ассоциативность
  k   нс/обращение
  1      0.232
  2      0.922
  3      0.921
  4      0.921
  5      0.921
  6      0.921
  7      0.922
  8      0.921
  9      2.580
 10      3.725
 11      4.667
 12      6.542
 13      6.222
 14      6.929
 15      6.639
 16      6.921
 17      6.742
 18      6.996
 19      6.866
 20      6.704
 21      6.629
 22      6.369
 23      6.655
 24      7.022
 25      6.696
 26      6.353
 27      7.618
 28      6.624
 29      7.305
 30      6.536
 31      6.546
 32      7.116
\end{lstlisting}
\par
Для $k$ от 2 до 8 ячеек время обращения постоянно и составляет около 0.92 нс. При $k = 9$ оно возрастает до 2.6 нс и далее растет до значения порядка 6.2-7.6 нс. Следовательно, в одном наборе кэша первого уровня помещаются 8 ячеек, то есть кэш \textbf{8-канальный множественно-ассоциативный} (8-way set associative). Значение при $k = 1$ (0.23 нс) в анализе не учитывается, так как цепочка из одной ячейки является особым случаем.

\newpage
